\section{Diseño de investigación}

El estudio integra dos diseños complementarios. Para el componente predictivo se utiliza un diseño longitudinal de series temporales con evaluación fuera de muestra. Los registros se ordenarán cronológicamente y se dividirán en entrenamiento, validación y prueba sin aleatorizar el tiempo. Todos los algoritmos utilizarán los mismos periodos y horizontes.

\[
\text{Datos válidos}\rightarrow\text{entrenamiento}\rightarrow
\text{validación temporal}\rightarrow\text{prueba futura}\rightarrow
\text{comparación}
\]

\begin{table}[H]
\caption{Diseño del componente predictivo}
\centering
\small
\begin{tabularx}{0.95\textwidth}{p{3.5cm}X}
\toprule
\textbf{Etapa} & \textbf{Aplicación en el estudio} \\
\midrule
Preparación & Cálculo del precio FOB unitario, agregación temporal, depuración de duplicados, tratamiento de nulos y construcción de rezagos. \\
Entrenamiento & Ajuste de Random Forest, HistGradientBoosting, ElasticNet y Prophet con el bloque histórico inicial. \\
Validación & Selección de hiperparámetros mediante ventanas temporales expansivas o deslizantes. \\
Prueba & Pronóstico de observaciones posteriores no utilizadas en el ajuste. \\
Comparación & Evaluación mediante MAE, RMSE, MAPE y $R^2$ sobre el mismo conjunto de prueba. \\
\bottomrule
\end{tabularx}
\end{table}

Para evaluar el apoyo a la decisión se aplicará un diseño preexperimental de un solo grupo:

\[
G\rightarrow O_1\rightarrow X\rightarrow O_2
\]

donde $G$ representa a los participantes de la prueba piloto; $O_1$, la medición inicial de comprensión, incertidumbre y forma de estimar rentabilidad; $X$, el uso guiado del prototipo; y $O_2$, la medición posterior. Este diseño permitirá estimar cambios iniciales de utilidad percibida, pero no autoriza atribuir causalidad definitiva por la ausencia de un grupo de control.
